\section{Degree-capped specializations}\label{sup:cap}
At four terms the elementary polynomial \(\dot x=2y,\ \dot y=-2x+\mu(-y+4y^3/3)\) has degree three. Its normalized first moment is proportional to \(E(E-1)\), so one small nonzero rational \(\mu\) gives a simple return. For one, two or three terms the zero bound is due to Buzzi, Carvalho and Gasull \cite[Theorem~1.2]{BuzziCarvalhoGasull2021}.

\subsection{Two simple roots at one focus and a two-sheeted lift}\label{sup:focus}
Use the \(k=2\) truncation of Gasull and Santana's (21) \cite[pp.~11--12]{GasullSantana2024}, as written in Section~\ref{sup:sources}, with \(b=a_2\). On \(y>0\), let \(u=\log y,\ w=x-1/y\). Then
\[
 \dot u=w,\quad \dot w=g(u)+f(u)w,\quad
 g=a_1(e^u-1)+b(e^{3u}-1),\quad f=a_1e^{2u}+e^{-u}.
\]
The point \((u,w)=(0,0)\) has \(A=-g'(0)=-a_1-3b>0\) near \((a_1,b)=(-1,0)\). Define \(V(u)=-\int_0^ug\), \(z=u\sqrt{2V(u)/(Au^2)}\). This real analytic chart has \(z'(0)=1\). Dividing the field by \(z'\) gives \(z'=w,\ w'=-Az+\widetilde f(z)w\), where primes on the field now refer to the new time, and \(\widetilde f(z)=f(u(z))\,du/dz\). Write \(c_{2j}=[z^{2j}]\widetilde f\). In particular \(c_0=1+a_1\). On \(a_1=-1\), differentiation of the primitive of \(f\) and Lagrange inversion give
\[
 c_{2j}=[u^{2j}](-e^{2u}+e^{-u})
 \left(\frac{2V(u)}{Au^2}\right)^{-(2j+1)/2},\quad
 V=e^u-1-u-b((e^{3u}-1)/3-u).
\]
\hypertarget{focus-jet-coefficients}{Expansion yields} \(c_2=9b/(3b-1)\), \(c_4|_{b=0}=1/3\). Thus \((c_0,c_2)\) are local coordinates in the two parameters. 

Let \(K_j=\int_0^{2\pi}\cos^{2j}\theta\sin^2\theta\,d\theta>0\) for \(j=0,1,2\), and \(C=K_2/3\). Choose small \(\delta>0\), with \(K_0c_0=4C\delta^4,\ K_1c_2=-5C\delta^2\). In \(z=\delta Z,\ w=\sqrt A\,\delta W\) on a compact annulus about radii 1 and 2, the time-normalized system has linear part \(Z'=W,\ W'=-Z\), with damping \(\widetilde f(\delta Z)/\sqrt A\). The odd part alone gives a center: the reflection \((Z,W,t)\mapsto(-Z,W,-t)\) fixes the vertical axis, and for small \(\delta\) every trajectory in the annulus crosses it twice with nonzero speed. Its one-turn displacement is identically zero. The even part divided by \(\delta^4\) tends in \(C^2\) to \(4C/K_0-5CZ^2/K_1+Z^4/3\). To see the leading return, multiply this bounded normalized part by an independent parameter \(\eta\), integrate the first energy variation on a section, and then set \(\eta=\delta^4\). Smooth dependence of solutions and their section derivatives on a compact annulus makes the \(O(\eta^2)\) remainder uniform in \(C^1\). As \(\delta\to0\) the odd centers converge to the circle; the change of the half squared radius, divided by \(\delta^4\), therefore tends, up to the chosen orientation, to
\begin{equation}\label{eq:capped-focus-return}
 C R^2(R^2-1)(R^2-4).
\end{equation}
Its two positive roots are simple. Equivalently one may take the explicit rational family \(a_1=-1+\delta^4/6,\ b=5\delta^2/54\) for all sufficiently small positive rational \(\delta\). Indeed \(K_0=\pi, K_1=\pi/4, K_2=\pi/8\), so the same two scaled coefficients have the required limits. Both cycles persist for a rational choice of \((a_1,b)\) near the selected parameters with \(a_1b(a_1+b)\ne0\). They lie in \(x,y>0\) near \((1,1)\). Set \(x=v^2\) alone and multiply the pullback by \(2v\). Its ordinary polynomial components are
\[
 P_*=a_1v^2 y^2+b y^3-(a_1+b),\qquad Q_*=2v^3y-2v.
\]
There are three and two distinct nonzero terms, of actual degree four. On each half-plane the pullback is a diffeomorphism and the speed factor never vanishes. On the time-reversed copy the section multiplier is the reciprocal of the original one. The two disjoint copies each contain two hyperbolic cycles. This one-coordinate specialization explains the four in the capped five-term row.

For the four-term truncation take \(b=0\) and \(a_1=-1-\delta^4/24\). Then \(K_0c_0=-C\delta^4\) and \(c_2=O(\delta^4)\), since \(c_2(-1,0)=0\). The same scaled return has leading polynomial \(CR^2(R^4-1)\), with one simple positive root. Squaring both coordinates as in Section~\ref{sup:sources} lifts this cycle to four hyperbolic cycles, one per quadrant, at degree seven with support \(\splitcount{2}{2}\). Applied to the five-term truncation with two simple returns from \eqref{eq:capped-focus-return}, the same map gives eight hyperbolic cycles, two per quadrant, at degree seven with support \(\splitcount{3}{2}\).

\paragraph{Three simple cycles in the six-term focus family.}
For \(k=3\) keep the additional \(c(e^{4u}-1)\) in \(g\), with \(b=a_2,\ c=a_3\); the damping \(f\) is unchanged. On \(a_1=-1\), \(A=1-3b-4c\). In the coefficient formula above replace \(V\) by \(V-c((e^{4u}-1)/4-u)\). Expansion gives
\[
 c_2=\frac{9(b+2c)}{3b+4c-1},\qquad
 c_4\big|_{b=-2c}=-\frac{19c-1}{3(2c+1)} .
\]
At \((a_1,b,c)=(-1,-2/19,1/19)\), \(A=21/19\), \(c_6=-4/21\), and
\(\det\partial(c_2,c_4)/\partial(b,c)=6859/147\). Thus \(c_0,c_2,c_4\) are local parameter coordinates. For \(j=0,1,2\) set
\[
 K_j c_{2j}=-(4/21)K_3\,\delta^{6-2j}
 [t^j]\big((t-1)(t-4)(t-9)\big).
\]
On compact annuli in the scaled radius \(R\), the odd damping again gives a center. After division by \(\delta^6\) the even damping converges in \(C^1\) to its prescribed polynomial. The first variation about the odd center tends to its circular average, as above. The normalized energy displacement has limit proportional to \(R^2(R^2-1)(R^2-4)(R^2-9)\), with nonzero proportionality constant. All three positive roots are simple. Rational approximation inside the resulting open parameter set preserves them and all four first-component coefficients. For \(k=1,2,3\), a single square \(x=v^2\) gives \((P(v^2,y),2vQ(v^2,y))\); the two squares give the formula in Section~\ref{sup:sources}. Each sheet carries the same simple fixed points, with inverse return map on a reversed sheet. The components in all three presentations are coprime: on the torus, \(xy=1\) restricts \(P\) to \(a_1y+\sum_{j=2}^k a_j y^{j+1}-a\), a nonzero polynomial. After the square substitutions the same identity holds on each lifted component. The constant term \(-a\ne0\) and the distinct coordinate factors in the lifted components exclude common axes.

\subsection{A quartic vertical potential at six, seven and eight terms}\label{sup:quartic}
Let \(V=(y^2-1)^2,\ U=A^2T_n(x)^2,\ (\dot x,\dot y)=(V',-U'+\mu F)\), and take \(F\) odd. On each of the \(n\) horizontal inverse branches of \(X=AT_n(x)\) near its simple zeros, after multiplication by \(A\) the absolute horizontal Jacobian tends in \(C^2\) to a positive constant on compact \(X\)-intervals. This also holds for \(n=1\), when the coordinate is linear. For
\[
 L_j(y)=\sum_{k=0}^j\binom{2j+1}{2k}y^{2j+1-2k}(2-y^2)^k
\]
and \(w=y^2-1\), the difference of the drive primitive \(L_j\) at \(a=\sqrt{1+w},b=\sqrt{1-w}\) is \((a-b)^{2j+1}\). Indeed \(L_j(a)=((a+b)^{2j+1}+(a-b)^{2j+1})/2\). The negative well has the same density integral by parity. Thus, up to a common positive normalization, the first moment of \(L_j\) in the disk \(X^2+w^2<E\) begins with \(k_j E^{j+1}\), where
\[
 k_j=(2j+1)\binom{2j}{j}/(4^j(j+1)).
\]
The identity follows by differentiating the odd trace and integrating the resulting \(w^{2j}\) on a disk. For fixed large \(A\), the even part of its analytic horizontal weight equals its value at zero plus \(O(X^2/A^2)\). The leading power and normalized coefficient are unchanged. Set
\[
 F_\delta=\sum_{j=0}^{d-1}b_j L_j,\qquad
 b_j=\frac{k_{d-1}}{k_j}\delta^{d-1-j}[s^j]\prod_{i=1}^{d-1}(s-i).
\]
Now \(b_{d-1}=1\), and on \(E=\delta s\) the sum of moments divided by \(\delta^d\) converges in \(C^1\) to \(k_{d-1}s\prod_{i=1}^{d-1}(s-i)\). The analytic expansion with the scaled unit disk justifies differentiation. This yields \(2n(d-1)\) simple inner roots.

For \(E>1\) an exterior oval surrounds both vertical wells. At constant horizontal weight and \(F=L_{d-1}\) its moment has the sign, up to one fixed orientation, of
\[
 J_d(E)=\int_0^{\pi/2}L_{d-1}'(\sqrt v)\cos^2\theta\sqrt{\sqrt E-1+v}\dd\theta,\qquad
 v=(1+\sqrt E)\sin^2\theta .
\]
This follows from \(y=\sqrt{1+\sqrt E}\sin\theta\) in the area integral. \hypertarget{quartic-moment-signs}{Table~\ref{tab:capped-exterior-signs} gives rational enclosures} for \(1000 J_d(E)\). The program \path{verify_low.py} reproduces these endpoints in \path{outputs/low-N.json}.
\begin{table}[htbp]\centering\small
\caption{Rational enclosures for \(1000J_d(E)\) at the exterior test energies.}\label{tab:capped-exterior-signs}
\begin{tabular}{rrrr}\toprule
\(d\)&\(E\)&lower&upper\\\midrule\inputrows{tex/quartic-rows.tex}\bottomrule
\end{tabular}
\end{table}
Divide \([0,\pi/2]\) into 8192 equal panels, enclose sine and cosine on the whole panel in 120-bit Arb arithmetic, evaluate the exact polynomial \(L'_{d-1}(\sqrt v)\) by Horner's rule in \(v\), and sum interval ranges times panel widths. This enclosure requires no quadrature remainder assumption. First take one large integer \(A\) retaining all strict exterior signs on all strips. Then take one small rational \(\delta>0\) retaining the signs and the inner simple brackets. On every exterior annulus adding \(\gamma y\) to \(F_\delta\) changes the first integral by \(\gamma\) times a strictly positive disk area. A common small rational regular value of the quotient, as in Section~\ref{sup:sources}, gives simple exterior roots while the inner roots persist. If an ordinary odd coefficient vanishes, add a sufficiently small nonzero rational in that vacant odd slot. The leading coefficient of \(L_{d-1}\) is \(\Re(1+i)^{2d-1}\ne0\). This fills exactly the \(d\) odd powers of \(y\), preserving all simple roots through the finite perturbations. One final small rational \(\mu>0\) realizes the cycles by the main return lemma. Their collars are disjoint within each strip and across strips. The support is \(\splitcount{2}{n+d}\), the actual degree is \(\max(3,2n-1,2d-1)\). For \((n,d)=(1,2),(2,2),(2,3),(2,4)\) this proves \(3,6,10,16\) cycles at degrees \(3,3,5,7\).


For the remaining stationary entries of Table~\othernumber{tab:small}{3} in the main article, Proposition~\othernumber{prop:stationary-budget}{6.1} there uses \((n,d)=(3,4),(4,4)\), giving \(25,35\) hyperbolic cycles at degree eight.
